\documentclass[aspectratio=169,10pt]{beamer}

% Build from this directory with a standard LaTeX toolchain, for example:
%   latexmk -pdf defense_presentation.tex

\usetheme{Madrid}
\usecolortheme{default}
\usefonttheme{professionalfonts}

\usepackage[T1]{fontenc}
\usepackage[utf8]{inputenc}
\usepackage{amsmath,amssymb,mathtools}
\usepackage{booktabs}
\usepackage{tabularx}
\usepackage{array}
\usepackage{graphicx}
\usepackage{tikz}
\usetikzlibrary{arrows.meta,positioning,shapes.geometric,fit}

\graphicspath{{assets/}}

\definecolor{bthblue}{RGB}{0,64,128}
\definecolor{bthcyan}{RGB}{0,150,190}
\definecolor{bthgray}{RGB}{65,70,75}
\definecolor{lightpanel}{RGB}{244,247,250}

\setbeamercolor{palette primary}{bg=bthblue,fg=white}
\setbeamercolor{palette secondary}{bg=bthcyan,fg=white}
\setbeamercolor{frametitle}{bg=bthblue,fg=white}
\setbeamercolor{title}{fg=bthblue}
\setbeamercolor{structure}{fg=bthblue}
\setbeamercolor{block title}{bg=bthblue,fg=white}
\setbeamercolor{block body}{bg=lightpanel,fg=black}
\setbeamertemplate{navigation symbols}{}
\setbeamertemplate{itemize item}{\raisebox{0.15ex}{\scriptsize$\blacktriangleright$}}
\setbeamertemplate{itemize subitem}{\raisebox{0.15ex}{\tiny$\blacktriangleright$}}

\newcommand{\alg}[1]{\textsc{#1}}
\newcommand{\R}{\mathbf{R}}
\newcommand{\C}{\mathbf{C}}
\newcommand{\Lmat}{\mathbf{L}}
\newcommand{\cp}{\mathbf{c}^{p}}
\newcommand{\cr}{\mathbf{c}^{r}}
\newcommand{\x}{\mathbf{x}}
\newcommand{\zt}{\mathbf{z}_{t}}
\newcommand{\lt}{\mathbf{l}_{t}}
\newcommand{\mt}{\mathbf{m}}
\newcommand{\yt}{\mathbf{y}}

\title[Private cloud scheduling]{Procurement planning and offline job scheduling in private clouds}
\subtitle{Thesis defense presentation}
\author[G. Karlsson]{Gustav Karlsson}
\institute[BTH]{Blekinge Institute of Technology\\Faculty of Computing}
\date{May 2026}

\begin{document}

{
\setbeamertemplate{footline}{}
\begin{frame}
  \begin{columns}[T,totalwidth=\textwidth]
    \begin{column}{0.72\textwidth}
      \vspace{0.5cm}
      {\usebeamerfont{title}\usebeamercolor[fg]{title}\LARGE\inserttitle\par}
      \vspace{0.35cm}
      {\large\insertsubtitle\par}
      \vspace{0.8cm}
      {\large\insertauthor\par}
      \vspace{0.15cm}
      {\normalsize Supervisor: Professor Lars Lundberg\par}
      \vspace{0.65cm}
      {\normalsize \insertinstitute\par}
      \vspace{0.15cm}
      {\normalsize \insertdate\par}
    \end{column}
    \begin{column}{0.22\textwidth}
      \raggedleft
      
\includegraphics[width=0.88\linewidth]{bthnotext.pdf}
    \end{column}
  \end{columns}
\end{frame}
}

\begin{frame}{Defense Roadmap}
  \begin{columns}[T,totalwidth=\textwidth]
    \begin{column}{0.47\textwidth}
      \begin{enumerate}
        \item Introduction and motivation
        \item Theory: bin-packing
        \item Problem description and model
        \item Related work and design gap
        \item Method: evaluated algorithms
      \end{enumerate}
    \end{column}
    \begin{column}{0.47\textwidth}
      \begin{enumerate}
        \setcounter{enumi}{5}
        \item Experimental methodology
        \item Results: cost and runtime
        \item Analysis and validity
        \item Conclusion and future work
      \end{enumerate}
    \end{column}
  \end{columns}
  \vspace{0.35cm}
  \begin{block}{Core thesis question}
    How can a private-cloud operator choose a hardware fleet and schedule predicted offline workloads at low total cost?
  \end{block}
\end{frame}

\section{Introduction}

\begin{frame}{Motivation: Private Clouds Need Planning}
  \begin{columns}[T,totalwidth=\textwidth]
    \begin{column}{0.50\textwidth}
      \begin{itemize}
        \item Organizations may avoid public cloud for regulatory, privacy, risk, or control reasons.
        \item In private clouds, the operator owns both the procurement decision and the scheduling decision.
        \item Energy-aware scheduling is not only about using fewer machines; it is about buying and powering the right machines over time.
      \end{itemize}
    \end{column}
    \begin{column}{0.46\textwidth}
      \begin{tikzpicture}[
        node distance=0.75cm,
        box/.style={rounded corners=2pt,draw=bthblue,thick,fill=lightpanel,align=center,minimum height=0.75cm,text width=3.0cm},
        arrow/.style={-{Latex[length=2mm]},thick,bthblue}
      ]
        \node[box] (predict) {Predicted workload};
        \node[box,below=of predict] (procure) {Hardware fleet\\procurement};
        \node[box,below=of procure] (schedule) {Offline job\\scheduling};
        \node[box,below=of schedule] (cost) {Total cost\\purchase + running};
        \draw[arrow] (predict) -- (procure);
        \draw[arrow] (procure) -- (schedule);
        \draw[arrow] (schedule) -- (cost);
      \end{tikzpicture}
    \end{column}
  \end{columns}
\end{frame}

\begin{frame}{Research Questions}
  \begin{block}{RQ1}
    How can we create a cost-efficient algorithm for offline multi-resource cloud procurement planning and job scheduling?
  \end{block}
  \begin{block}{RQ2}
    How can we create a procurement planner and cloud job scheduler which is optimized for both scheduling quality and execution time?
  \end{block}
  \vspace{0.25cm}
  \begin{itemize}
    \item The thesis treats complete workload demand as known for a finite planning horizon.
    \item Each job must run in its assigned time slot and cannot be split across machines.
  \end{itemize}
\end{frame}

\section{Theory}

\begin{frame}{Bin-Packing View}
  \begin{columns}[T,totalwidth=\textwidth]
    \begin{column}{0.47\textwidth}
      \begin{block}{Classical bin-packing}
        Pack items into bins while respecting capacity constraints and minimizing the number or cost of opened bins.
      \end{block}
      \begin{itemize}
        \item Jobs become multi-resource items.
        \item Machines become heterogeneous bins.
        \item Opening a bin has purchase and running cost consequences.
      \end{itemize}
    \end{column}
    \begin{column}{0.49\textwidth}
      \centering
      \begin{tikzpicture}[
        item/.style={draw=bthcyan,fill=bthcyan!15,rounded corners=1pt,minimum width=0.75cm,minimum height=0.35cm},
        bin/.style={draw=bthblue,thick,minimum width=2.4cm,minimum height=1.55cm},
        label/.style={font=\scriptsize}
      ]
        \node[item] at (0,1.2) {};
        \node[item] at (0.9,1.2) {};
        \node[item] at (1.8,1.2) {};
        \node[label] at (0.9,0.73) {jobs};
        \node[bin] at (0, -0.35) {};
        \node[bin] at (2.9, -0.35) {};
        \node[label] at (0,-1.35) {machine type 1};
        \node[label] at (2.9,-1.35) {machine type 2};
        \node[item,minimum width=0.65cm] at (-0.45,-0.55) {};
        \node[item,minimum width=0.95cm] at (0.35,-0.12) {};
        \node[item,minimum width=0.55cm] at (2.45,-0.48) {};
        \draw[-{Latex[length=2mm]},thick,bthgray] (0.9,0.95) -- (0.2,0.35);
        \draw[-{Latex[length=2mm]},thick,bthgray] (1.3,0.95) -- (2.6,0.35);
      \end{tikzpicture}
    \end{column}
  \end{columns}
  \vspace{0.15cm}
  \[
    \sum_i \mathbf{s}(i)x_{ij} \leq \mathbf{c}(j)y_j
  \]
\end{frame}

\begin{frame}{Problem Inputs}
  \begin{columns}[T,totalwidth=\textwidth]
    \begin{column}{0.48\textwidth}
      \begin{align*}
        \R &= \begin{bmatrix}
          | & | & & |\\
          \mathbf{r}_1 & \mathbf{r}_2 & \cdots & \mathbf{r}_J\\
          | & | & & |
        \end{bmatrix}
        &&\in \mathbb{Z}_{\geq 0}^{K \times J}\\[0.3cm]
        \C &= \begin{bmatrix}
          | & | & & |\\
          \mathbf{m}_1 & \mathbf{m}_2 & \cdots & \mathbf{m}_M\\
          | & | & & |
        \end{bmatrix}
        &&\in \mathbb{Z}_{\geq 0}^{K \times M}
      \end{align*}
    \end{column}
    \begin{column}{0.48\textwidth}
      \begin{align*}
        \Lmat &= \begin{bmatrix}
          | & | & & |\\
          \mathbf{l}_1 & \mathbf{l}_2 & \cdots & \mathbf{l}_T\\
          | & | & & |
        \end{bmatrix}
        &&\in \mathbb{Z}_{\geq 0}^{J \times T}\\[0.3cm]
        \cp &= \C^\top \boldsymbol{\alpha}, \qquad
        \cr = \gamma \cp
      \end{align*}
    \end{column}
  \end{columns}
  \vspace{0.2cm}
  \begin{itemize}
    \item $K$: resources; $J$: job types; $M$: machine types; $T$: time slots.
    \item $\boldsymbol{\alpha}$ weights resources by cost; $\gamma$ maps purchase cost to per-slot running cost.
  \end{itemize}
\end{frame}

\begin{frame}{Lower-Bound Procurement Model}
  \begin{block}{Simplified model}
    This model checks aggregate capacity, but does not guarantee that individual jobs can be packed into individual machines.
  \end{block}
  \begin{align*}
    \min_{\x,\zt}\quad
      & \x^\top\cp + \sum_{t=1}^{T}\sum_{i=1}^{M} z_{t,i}c_i^r\\
    \text{s.t.}\quad
      & \C\zt \geq \R\lt, && \forall t\\
      & \zt \leq \x, && \forall t\\
      & \x \in \mathbb{Z}_{\geq 0}^{M},\quad
        \zt \in \mathbb{Z}_{\geq 0}^{M}, && \forall t
  \end{align*}
  \vspace{-0.1cm}
  \begin{itemize}
    \item Useful as a lower bound, but insufficient for valid schedules.
  \end{itemize}
\end{frame}

\begin{frame}{Packing-Aware Model}
  \begin{columns}[T,totalwidth=\textwidth]
    \begin{column}{0.48\textwidth}
      \begin{block}{Valid packing configurations}
        \[
          S_i = \{\mathbf{y}\in\mathbb{Z}_{\geq 0}^{J}\mid \R\mathbf{y}\leq\mathbf{m}_i\}
        \]
        Each element of $S_i$ is a feasible set of jobs for one machine of type $i$.
      \end{block}
      \[
        \mathbf{Y}_i =
        \begin{bmatrix}
          | & | & & |\\
          \mathbf{y}_{i,1} & \mathbf{y}_{i,2} & \cdots & \mathbf{y}_{i,|S_i|}\\
          | & | & & |
        \end{bmatrix}
      \]
    \end{column}
    \begin{column}{0.48\textwidth}
      \begin{align*}
        \min\quad
          & \x^\top\cp
          + \sum_{i=1}^{M} c_i^r
            \sum_{t=1}^{T}\sum_{k=1}^{|S_i|} n_{i,t,k}\\
        \text{s.t.}\quad
          & \sum_{i=1}^{M}\mathbf{Y}_i\mathbf{n}_{i,t} \geq \lt,
          && \forall t\\
          & \sum_{k=1}^{|S_i|} n_{i,t,k} \leq x_i,
          && \forall i,t\\
          & \x,\mathbf{n}_{i,t}\in\mathbb{Z}_{\geq 0}
      \end{align*}
    \end{column}
  \end{columns}
\end{frame}

\section{Method}

\begin{frame}{Algorithms Evaluated}
  \begin{table}
    \centering
    \scriptsize
    \begin{tabularx}{\textwidth}{@{}lX@{}}
      \toprule
      Algorithm & Sorting and placement idea\\
      \midrule
      \alg{FFD} / \alg{FFDLex} & First-fit decreasing with lexicographic ordering.\\
      \alg{FFDSum} & Sort jobs by decreasing sum of resource demands.\\
      \alg{FFDProd} & Sort jobs by decreasing product of resource demands.\\
      \alg{FFDMax} & Sort jobs by largest single resource demand.\\
      \alg{FFDL2} & Sort jobs by Euclidean norm of resource demand vector.\\
      \alg{FFDNew} & Weighted ordering plus slack-aware new-machine selection.\\
      \alg{BFD} & Resource-weighted, cost-aware best-fit over open and new machines.\\
      \bottomrule
    \end{tabularx}
  \end{table}
  \vspace{0.2cm}
  \begin{itemize}
    \item All algorithms select feasible machine types and produce validated schedules.
    \item The central comparison is quality versus runtime overhead.
  \end{itemize}
\end{frame}

\begin{frame}{BFD and FFDNew: Shared Setup}
  \begin{columns}[T,totalwidth=\textwidth]
    \begin{column}{0.50\textwidth}
      \begin{enumerate}
        \item For each time slot, start with no open machines.
        \item Compute weighted job priority:
        \[
          \mathbf{v}=\R^\top\boldsymbol{\alpha}
        \]
        \item Sort job types by decreasing $v_j$.
        \item Pack all remaining jobs of the current type before moving to the next type.
      \end{enumerate}
    \end{column}
    \begin{column}{0.46\textwidth}
      \begin{block}{Common placement count}
        \[
          q=\min_{k:\hat{r}_{j,k}>0}
            \left\lfloor\frac{z_k}{\hat{r}_{j,k}}\right\rfloor
        \]
        \[
          n=\min(q,\eta),\qquad
          \mathbf{u}=\mathbf{z}-n\hat{\mathbf{r}}_j
        \]
        $q$ is how many jobs fit in the candidate capacity vector $\mathbf{z}$; $n$ is how many are actually placed.
      \end{block}
    \end{column}
  \end{columns}
\end{frame}

\begin{frame}{BFD: Open-Bin Score \texorpdfstring{$\Phi_b$}{Phi}}
  \begin{columns}[T,totalwidth=\textwidth]
    \begin{column}{0.48\textwidth}
      \begin{block}{Remaining capacity of open bin $b$}
        If $i=\tau_{t,b}$ is the machine type of open bin $b$:
        \[
          \boldsymbol{\rho}_b
          =
          \mathbf{m}_i
          -
          \sum_{j'=1}^{J} y_{t,j',b}\hat{\mathbf{r}}_{j'}
        \]
      \end{block}
      \begin{itemize}
        \item For open bins, the candidate capacity is $\mathbf{z}=\boldsymbol{\rho}_b$.
        \item Only bins with $q_b\geq 1$ are feasible.
      \end{itemize}
    \end{column}
    \begin{column}{0.48\textwidth}
      \begin{block}{BFD selection key}
        \[
          \Phi_b
          =
          \sum_{k=1}^{K}
          \alpha_k
          \left(\rho_{b,k}-n_b\hat{r}_{j,k}\right)^2
        \]
        \[
          \kappa_b =
          \left(\Phi_b,\;c^r_{\tau_{t,b}},\;b\right)
        \]
      \end{block}
      \begin{itemize}
        \item Select $b^*=\arg\,\min_{b:q_b\geq 1}\kappa_b$.
        \item The score favors tight fits in expensive or scarce resource dimensions.
      \end{itemize}
    \end{column}
  \end{columns}
\end{frame}

\begin{frame}{New-Machine Score \texorpdfstring{$\Psi_i$}{Psi}: Shared by BFD and FFDNew}
  \begin{columns}[T,totalwidth=\textwidth]
    \begin{column}{0.48\textwidth}
      \begin{block}{Empty machine type candidate}
        When no open bin is feasible, evaluate every machine type $i$:
        \[
          q_i=
          \min_{k:\hat{r}_{j,k}>0}
          \left\lfloor
            \frac{m_{i,k}}{\hat{r}_{j,k}}
          \right\rfloor
        \]
        \[
          n_i=\min(\eta,q_i)
        \]
      \end{block}
      \begin{itemize}
        \item Candidate capacity is the full machine vector $\mathbf{z}=\mathbf{m}_i$.
        \item Only machine types with $q_i\geq 1$ can be opened.
      \end{itemize}
    \end{column}
    \begin{column}{0.48\textwidth}
      \begin{block}{Normalized new-bin score}
        \[
          \Psi_i
          =
          \frac{
            \sum_{k=1}^{K}
            \alpha_k
            \left(m_{i,k}-n_i\hat{r}_{j,k}\right)^2
          }{c_i^p}
        \]
        \[
          \kappa_i =
          \left(\Psi_i,\;c_i^p+c_i^r,\;i\right)
        \]
      \end{block}
      \begin{itemize}
        \item Select $i^*=\arg\,\min_{i:q_i\geq 1}\kappa_i$.
        \item Dividing by $c_i^p$ avoids unfairly penalizing larger machine types only because they have more raw capacity.
      \end{itemize}
    \end{column}
  \end{columns}
\end{frame}

\begin{frame}{FFDNew: Same New-Bin Logic, Simpler Open-Bin Logic}
  \begin{columns}[T,totalwidth=\textwidth]
    \begin{column}{0.50\textwidth}
      \begin{block}{BFD}
        \begin{enumerate}
          \item Score all feasible open bins with $\Phi_b$.
          \item Choose the minimum key $\kappa_b$.
          \item If none fit, open a new machine type using $\Psi_i$.
        \end{enumerate}
      \end{block}
    \end{column}
    \begin{column}{0.46\textwidth}
      \begin{block}{FFDNew}
        \begin{enumerate}
          \item Scan open bins in the order they were opened.
          \item Select the first bin with $q_b\geq 1$:
          \[
            b^*=\arg\,\min_{b\in B:q_b\geq 1} b
          \]
          \item If none fit, open a new machine type using the same $\Psi_i$ as BFD.
        \end{enumerate}
      \end{block}
    \end{column}
  \end{columns}
  \vspace{0.15cm}
  \begin{itemize}
    \item FFDNew is therefore a hybrid: BFD-style ordering and new-machine selection, but first-fit selection among already-open bins.
  \end{itemize}
\end{frame}

\begin{frame}{Experimental Datasets}
  \begin{table}
    \centering
    \small
    \begin{tabular}{@{}lccccc@{}}
      \toprule
      Dataset & Instances & $K$ & $J$ range & $M$ range & $T$ range\\
      \midrule
      Balanced & 100 & 4 & 6--8 & 6--8 & 100--200\\
      Job-heavy & 100 & 4 & 12--16 & 6--8 & 100--200\\
      Machine-heavy & 100 & 4 & 6--8 & 12--16 & 100--200\\
      \bottomrule
    \end{tabular}
  \end{table}
  \vspace{0.25cm}
  \begin{itemize}
    \item Resource types: CPU, memory, disk, and I/O.
    \item Machine, job, and time-slot generation uses controlled specialization rather than uniform sampling.
    \item Each dataset uses 100 independently generated problem instances.
  \end{itemize}
\end{frame}

\begin{frame}{Instance Generation Structure}
  \begin{columns}[T,totalwidth=\textwidth]
    \begin{column}{0.50\textwidth}
      \begin{block}{Baseline plus jitter}
        \[
          C_{k,i}\leftarrow c_{0,k}U([c_{\min},c_{\max}])
        \]
        \[
          R_{k,j}\leftarrow d_{0,k}U([d_{\min},d_{\max}])
        \]
      \end{block}
    \end{column}
    \begin{column}{0.46\textwidth}
      \begin{block}{Correlated specialization}
        \[
          \mathbf{q}^{\mathrm{machine}}
          =
          \eta^{\mathrm{machine}}\mathbf{q}^{\mathrm{job}}
          +(1-\eta^{\mathrm{machine}})\frac{\mathbf{1}_K}{K}
        \]
      \end{block}
    \end{column}
  \end{columns}
  \vspace{0.2cm}
  \begin{itemize}
    \item Some jobs are CPU-, memory-, disk-, or I/O-heavy.
    \item Some machines are specialized toward matching resource dimensions.
    \item This creates structured heterogeneity similar to private-cloud fleet planning.
  \end{itemize}
\end{frame}

\begin{frame}{Problem Instance Generation Pipeline}
  \centering
  \begin{tikzpicture}[
    scale=0.88,
    transform shape,
    node distance=0.85cm,
    box/.style={rounded corners=2pt,draw=bthblue,thick,fill=lightpanel,align=center,minimum height=0.72cm,text width=3.2cm},
    smallbox/.style={rounded corners=2pt,draw=bthcyan,thick,fill=bthcyan!10,align=center,minimum height=0.64cm,text width=2.7cm},
    arrow/.style={-{Latex[length=2mm]},thick,bthblue}
  ]
    \node[box] (dims) {Sample dimensions\\$K,M,J,T$};
    \node[box,right=of dims] (jobp) {Assign job primary resources\\$\mathbf{p}^{\mathrm{job}}$};
    \node[box,right=of jobp] (hist) {Build job-resource histogram\\$\mathbf{q}^{\mathrm{job}}$};
    \node[smallbox,below=of jobp] (R) {Generate job demands\\$\R$};
    \node[smallbox,below=of hist] (C) {Generate machine capacities\\$\C$};
    \node[smallbox,right=of hist] (L) {Generate time slots\\$\Lmat$};
    \node[box,below=of C] (cost) {Compute costs\\$\cp,\cr,\boldsymbol{\alpha}$};
    \draw[arrow] (dims) -- (jobp);
    \draw[arrow] (jobp) -- (hist);
    \draw[arrow] (jobp) -- (R);
    \draw[arrow] (hist) -- (C);
    \draw[arrow] (hist) -- (L);
    \draw[arrow] (R) -- (C);
    \draw[arrow] (C) -- (cost);
  \end{tikzpicture}
  \vspace{0.25cm}
  \begin{itemize}
    \item The output is random, but not independent noise: supply, demand, and temporal load are coupled through primary resources.
    \item The generated instance is $(\C,\R,\Lmat,\cp,\cr,\boldsymbol{\alpha})$.
  \end{itemize}
\end{frame}

\begin{frame}{Primary Resources: Creating Structured Heterogeneity}
  \begin{columns}[T,totalwidth=\textwidth]
    \begin{column}{0.48\textwidth}
      \begin{block}{Choose primary resources}
        For $n$ objects and specialization rate $\rho$:
        \[
          s=\mathrm{round}(n\rho)
        \]
        Select $s$ objects without replacement, then assign each selected object a resource
        \[
          k\sim\mathrm{Categorical}(\mathbf{q})
        \]
      \end{block}
      \begin{itemize}
        \item Unselected objects have no primary resource.
        \item Each selected object has at most one primary resource.
      \end{itemize}
    \end{column}
    \begin{column}{0.48\textwidth}
      \begin{block}{Correlation from job profile}
        Job types are sampled first. Their primary-resource histogram becomes:
        \[
          \mathbf{q}^{\mathrm{job}}
          =
          \frac{\mathbf{h}}{\sum_{k=1}^{K}h_k}
        \]
        Machine and slot probabilities blend this profile with uniform randomness:
        \[
          \mathbf{q}^{\mathrm{machine}}
          =
          \eta^{\mathrm{machine}}\mathbf{q}^{\mathrm{job}}
          +(1-\eta^{\mathrm{machine}})\frac{\mathbf{1}_K}{K}
        \]
      \end{block}
    \end{column}
  \end{columns}
\end{frame}

\begin{frame}{Generating Job and Machine Types}
  \begin{columns}[T,totalwidth=\textwidth]
    \begin{column}{0.48\textwidth}
      \begin{block}{Job demand matrix $\R$}
        Initialize every demand with a base value and jitter:
        \[
          R_{k,j}\leftarrow d_{0,k}U([d_{\min},d_{\max}])
        \]
        If job type $j$ has primary resource $k^*$:
        \[
          R_{k^*,j}\leftarrow uR_{k^*,j},
          \qquad
          u\sim\mathrm{UniformInteger}([u_{\min},u_{\max}])
        \]
      \end{block}
    \end{column}
    \begin{column}{0.48\textwidth}
      \begin{block}{Machine capacity matrix $\C$}
        Initialize every capacity similarly:
        \[
          C_{k,i}\leftarrow c_{0,k}U([c_{\min},c_{\max}])
        \]
        If machine type $i$ has primary resource $k^*$:
        \[
          C_{k^*,i}\leftarrow uC_{k^*,i}
        \]
      \end{block}
    \end{column}
  \end{columns}
  \vspace{0.15cm}
  \begin{itemize}
    \item Values are rounded and clamped to at least 1.
    \item If any job type fits on no machine, capacity is added to a selected machine type, preferably one with a matching primary resource.
  \end{itemize}
\end{frame}

\begin{frame}{Generating Time Slots and Workload Matrix}
  \begin{columns}[T,totalwidth=\textwidth]
    \begin{column}{0.48\textwidth}
      \begin{block}{Slot primary resources}
        Time-slot specialization follows the job-type profile:
        \[
          \mathbf{q}^{\mathrm{slot}}
          =
          \eta^{\mathrm{slot}}\mathbf{q}^{\mathrm{job}}
          +(1-\eta^{\mathrm{slot}})\frac{\mathbf{1}_K}{K}
        \]
        Then:
        \[
          \mathbf{p}^{\mathrm{slot}}
          =
          \mathrm{ChoosePrimaryResources}(T,\rho^{\mathrm{slot}},\mathbf{q}^{\mathrm{slot}})
        \]
      \end{block}
    \end{column}
    \begin{column}{0.48\textwidth}
      \begin{block}{Slot job counts}
        Total jobs in slot $t$:
        \[
          N_t=\max(1,\mathrm{round}(\lambda_0 u)),
          \qquad
          u\sim U([\lambda_{\min},\lambda_{\max}))
        \]
        After emphasizing matching job types:
        \[
          \boldsymbol{\ell}_t
          \sim
          \mathrm{Multinomial}(N_t,\boldsymbol{\phi})
        \]
      \end{block}
    \end{column}
  \end{columns}
  \vspace{0.15cm}
  \begin{itemize}
    \item A slot with primary resource $k^*$ increases the weights of job types that also have primary resource $k^*$.
    \item This creates temporal bursts of related job types while preserving random variation in total load.
  \end{itemize}
\end{frame}

\begin{frame}{From Problem Instances to Datasets}
  \begin{columns}[T,totalwidth=\textwidth]
    \begin{column}{0.48\textwidth}
      \begin{block}{Per-instance sampling}
        For each generated instance:
        \[
          J\sim U(I_J),\qquad
          M\sim U(I_M),\qquad
          T\sim U(I_T)
        \]
        with $K=4$ fixed to CPU, memory, disk, and I/O.
      \end{block}
    \end{column}
    \begin{column}{0.48\textwidth}
      \begin{block}{Three dataset regimes}
        \begin{itemize}
          \item Balanced: $J=6$--$8$, $M=6$--$8$.
          \item Job-heavy: $J=12$--$16$, $M=6$--$8$.
          \item Machine-heavy: $J=6$--$8$, $M=12$--$16$.
        \end{itemize}
      \end{block}
    \end{column}
  \end{columns}
  \vspace{0.15cm}
  \begin{itemize}
    \item Each dataset contains 100 generated problem instances using fixed seed 5000.
    \item The regimes isolate whether extra job diversity or extra machine diversity changes scheduler behavior.
  \end{itemize}
\end{frame}

\begin{frame}{Evaluation Metrics}
  \begin{columns}[T,totalwidth=\textwidth]
    \begin{column}{0.48\textwidth}
      \begin{block}{Cost}
        \[
          \text{total cost} =
          \x^\top\cp + \sum_{t=1}^{T}\zt^\top\cr
        \]
        Reported as average, min, max, and paired ratios.
      \end{block}
    \end{column}
    \begin{column}{0.48\textwidth}
      \begin{block}{Performance profile}
        \[
          \rho_s(\tau)=
          \frac{1}{|P|}
          \left|\{p\in P:r_{p,s}\leq\tau\}\right|
        \]
        Summarizes how often solver $s$ is within factor $\tau$ of the best solver.
      \end{block}
    \end{column}
  \end{columns}
  \vspace{0.2cm}
  \begin{itemize}
    \item Runtime is measured per problem instance.
    \item Paired two-tailed $t$-tests compare per-instance ratios.
  \end{itemize}
\end{frame}

\section{Results}

\begin{frame}{Cost Summary: Best Two Algorithms Dominate}
  \begin{table}
    \centering
    \scriptsize
    \begin{tabular}{@{}lrrrr@{}}
      \toprule
      Dataset / algorithm & Avg. cost & Avg. runtime & Avg. machines & Rank\\
      \midrule
      Balanced \alg{FFDNew} & 8132.64 & 0.07s & 32.26 & 1\\
      Balanced \alg{BFD} & 8134.48 & 0.07s & 32.26 & 2\\
      Balanced best simple FFD & 8943.51 & 0.05s & 24.63 & 3\\
      \addlinespace
      Job-heavy \alg{FFDNew} & 8433.94 & 0.08s & 34.89 & 1\\
      Job-heavy \alg{BFD} & 8437.84 & 0.07s & 34.88 & 2\\
      Job-heavy best simple FFD & 8991.99 & 0.05s & 25.61 & 3\\
      \addlinespace
      Machine-heavy \alg{FFDNew} & 6982.95 & 0.09s & 33.32 & 1\\
      Machine-heavy \alg{BFD} & 6983.41 & 0.09s & 33.32 & 2\\
      Machine-heavy best simple FFD & 8053.82 & 0.05s & 27.37 & 3\\
      \bottomrule
    \end{tabular}
  \end{table}
  \vspace{-0.1cm}
  \begin{itemize}
    \item \alg{BFD} and \alg{FFDNew} are consistently the two lowest-cost schedulers.
  \end{itemize}
\end{frame}

\begin{frame}{BFD vs. FFDNew: Paired Cost Ratios}
  \begin{table}
    \centering
    \small
    \begin{tabular}{@{}lrrrrl@{}}
      \toprule
      Dataset & $n$ & Mean ratio & 95\% CI & $p$-value & Decision\\
      \midrule
      Balanced & 100 & 1.00028 & 0.999877--1.00069 & 0.170 & Fail to reject\\
      Job-heavy & 100 & 1.00044 & 1.00002--1.00086 & 0.0407 & Reject\\
      Machine-heavy & 100 & 1.00011 & 0.999771--1.00045 & 0.522 & Fail to reject\\
      \bottomrule
    \end{tabular}
  \end{table}
  \vspace{0.2cm}
  \begin{block}{Interpretation}
    Even where the null hypothesis is rejected, the effect size is very small: the job-heavy mean ratio differs from 1 by about 0.044\%.
  \end{block}
\end{frame}

\begin{frame}{Performance Profiles: Balanced Dataset}
  \centering
  \includegraphics[width=0.88\textwidth,height=0.72\textheight,keepaspectratio]{eval_performance_profiles_balanced.pdf}
  \vspace{-0.1cm}
  \begin{itemize}
    \item \alg{FFDNew} and \alg{BFD} tie on 81 out of 100 balanced instances.
  \end{itemize}
\end{frame}

\begin{frame}{Performance Profiles: Job-Heavy and Machine-Heavy}
  \begin{columns}[T,totalwidth=\textwidth]
    \begin{column}{0.49\textwidth}
      \centering
      \includegraphics[width=\linewidth,height=0.62\textheight,keepaspectratio]{eval_performance_profiles_job_heavy.pdf}\\
      \scriptsize Job-heavy: \alg{FFDNew} and \alg{BFD} tie on 43 instances.
    \end{column}
    \begin{column}{0.49\textwidth}
      \centering
      \includegraphics[width=\linewidth,height=0.62\textheight,keepaspectratio]{eval_performance_profiles_machine_heavy.pdf}\\
      \scriptsize Machine-heavy: \alg{FFDNew} and \alg{BFD} tie on 86 instances.
    \end{column}
  \end{columns}
\end{frame}

\begin{frame}{Cost Ratio Distributions}
  \begin{columns}[T,totalwidth=\textwidth]
    \begin{column}{0.33\textwidth}
      \centering
      \includegraphics[width=\linewidth,height=0.46\textheight,keepaspectratio]{eval_cost_ratio_dist_balanced_bfd_vs_ffd_new.pdf}\\
      \scriptsize Balanced
    \end{column}
    \begin{column}{0.33\textwidth}
      \centering
      \includegraphics[width=\linewidth,height=0.46\textheight,keepaspectratio]{eval_cost_ratio_dist_job_heavy_bfd_vs_ffd_new.pdf}\\
      \scriptsize Job-heavy
    \end{column}
    \begin{column}{0.33\textwidth}
      \centering
      \includegraphics[width=\linewidth,height=0.46\textheight,keepaspectratio]{eval_cost_ratio_dist_machine_heavy_bfd_vs_ffd_new.pdf}\\
      \scriptsize Machine-heavy
    \end{column}
  \end{columns}
  \vspace{0.2cm}
  \begin{itemize}
    \item The distributions are tightly concentrated around a cost ratio of 1.
  \end{itemize}
\end{frame}

\begin{frame}{Runtime: Quality Comes With Small Overhead}
  \centering
  \includegraphics[width=0.86\textwidth,height=0.68\textheight,keepaspectratio]{eval_runtime_cross_dataset.pdf}
  \vspace{-0.05cm}
  \begin{itemize}
    \item All average runtimes remain below 100 ms per instance.
    \item \alg{BFD} and \alg{FFDNew} are slower than simpler FFD variants, but the absolute overhead is small.
  \end{itemize}
\end{frame}

\section{Analysis}

\begin{frame}{Why Lower Cost Can Use More Machines}
  \begin{table}
    \centering
    \small
    \begin{tabular}{@{}lrrr@{}}
      \toprule
      Dataset & Mean machine ratio & 95\% CI & $p$-value\\
      \midrule
      Balanced & 1.323 & 1.273--1.373 & $6.70\cdot 10^{-23}$\\
      Job-heavy & 1.376 & 1.330--1.422 & $1.34\cdot 10^{-29}$\\
      Machine-heavy & 1.245 & 1.193--1.297 & $3.06\cdot 10^{-15}$\\
      \bottomrule
    \end{tabular}
  \end{table}
  \vspace{0.25cm}
  \begin{block}{Interpretation}
    \alg{BFD} uses more machines than \alg{FFDMax}, yet achieves lower total cost. In this model, cost efficiency comes from choosing cheaper and better-fitting machine types, not merely minimizing active machine count.
  \end{block}
\end{frame}

\begin{frame}{Dataset-Level Interpretation}
  \begin{itemize}
    \item The machine-heavy dataset has more available machine types, increasing placement flexibility.
    \item For \alg{BFD}, machine-heavy instances had lower mean cost than balanced instances:
      \[
        6983 \quad \text{vs.}\quad 8134
      \]
    \item The paired one-tailed cross-dataset test supported this directional hypothesis:
      \[
        t=-2.279,\qquad p=0.01242,\qquad \Delta=-1151
      \]
    \item This is consistent with the generator: more machine types raise the chance of finding a machine specialized for the workload.
  \end{itemize}
\end{frame}

\begin{frame}{Answering the Research Questions}
  \begin{block}{RQ1: Cost-efficient offline scheduling}
    Traditional packing heuristics can be adapted to private-cloud procurement planning by adding heterogeneous machine types, purchase costs, running costs, and multi-resource feasibility checks.
  \end{block}
  \begin{block}{RQ2: Quality and execution time}
    \alg{BFD} and \alg{FFDNew} produce the lowest observed solution costs. Their runtimes are higher than simpler heuristics, but remain below 100 ms per instance in the evaluation.
  \end{block}
  \vspace{0.15cm}
  \begin{itemize}
    \item The practical trade-off is mild in the studied setting: the best quality-oriented algorithms remain fast.
  \end{itemize}
\end{frame}

\begin{frame}{Validity Considerations}
  \begin{columns}[T,totalwidth=\textwidth]
    \begin{column}{0.48\textwidth}
      \begin{block}{Internal validity}
        \begin{itemize}
          \item Random generator may favor some heuristics.
          \item Fixed seeds improve reproducibility but do not remove seed sensitivity.
          \item Pipeline bugs or scheduler implementation errors could affect results.
        \end{itemize}
      \end{block}
    \end{column}
    \begin{column}{0.48\textwidth}
      \begin{block}{External validity}
        \begin{itemize}
          \item Synthetic workloads omit bursts, seasonality, and production-specific correlations.
          \item Cost model is linear and simplified.
          \item Startup delays, maintenance windows, and multi-slot jobs are not modeled.
        \end{itemize}
      \end{block}
    \end{column}
  \end{columns}
\end{frame}

\section{Conclusion}

\begin{frame}{Conclusion}
  \begin{itemize}
    \item The thesis models private-cloud procurement and offline job scheduling as a cost-aware, multi-resource, heterogeneous bin-packing problem.
    \item \alg{BFD} and \alg{FFDNew} dominate the evaluated algorithms on solution cost.
    \item Their cost difference is practically very small across all datasets.
    \item Although they are slower than simpler FFD variants, all average runtimes remain below 100 ms per instance.
    \item Lower cost does not necessarily mean fewer machines; it can mean a more cost-efficient allocation across machine types.
  \end{itemize}
\end{frame}

\begin{frame}{Future Work}
  \begin{itemize}
    \item Treat workload predictions probabilistically rather than deterministically.
    \item Evaluate advanced search methods such as simulated annealing, tabu search, variable neighborhood search, or hybrid exact-heuristic approaches.
    \item Extend the objective beyond total cost, e.g. jointly minimize cost, machine count, or number of machine types used.
    \item Evaluate on realistic production traces such as the Google Borg cluster workload trace.
    \item Add operational constraints: startup delays, shutdown delays, maintenance windows, and jobs spanning multiple time slots.
    \item Study theoretical bounds and worst-case instances for the proposed algorithms.
  \end{itemize}
\end{frame}

{
\setbeamertemplate{footline}{}
\begin{frame}
  \begin{columns}[c,totalwidth=\textwidth]
    \begin{column}{0.70\textwidth}
      {\usebeamerfont{title}\usebeamercolor[fg]{title}\Huge Thank you\par}
      \vspace{0.45cm}
      {\Large Questions?\par}
      \vspace{0.8cm}
      {\normalsize Gustav Karlsson\\Blekinge Institute of Technology\par}
    \end{column}
    \begin{column}{0.22\textwidth}
      \raggedleft
      
\includegraphics[width=0.92\linewidth]{bthnotext.pdf}
    \end{column}
  \end{columns}
\end{frame}
}

\end{document}
